\documentclass[a4paper,11pt]{article}
\usepackage[utf8]{inputenc}
\usepackage[T1]{fontenc}
\usepackage[french]{babel}
\usepackage{amsmath,amssymb}
\usepackage{geometry}
\usepackage{tikz}
\usepackage{enumitem}
\usepackage{fancyhdr}
\usepackage{tcolorbox}
\usepackage{pgfplots}
\pgfplotsset{compat=1.18}

\geometry{margin=2cm}
\setlength{\headheight}{14pt}

\pagestyle{fancy}
\fancyhf{}
\lhead{BTS -- Mathématiques}
\rhead{Évaluation -- Séries de Fourier}
\cfoot{\thepage}

\tcbset{
  exobox/.style={colback=orange!5!white, colframe=orange!60!black, fonttitle=\bfseries, title=#1},
  reponsebox/.style={colback=gray!5!white, colframe=gray!40!black, sharp corners, boxrule=0.5pt, left=4pt, right=4pt, top=4pt, bottom=4pt}
}

\begin{document}

\begin{center}
  {\LARGE \textbf{Évaluation -- Séries de Fourier}} \\[0.2cm]
  {\Large Calcul de $a_0$ : valeur moyenne d'un signal périodique} \\[0.2cm]
  \rule{0.6\textwidth}{0.4pt} \\[0.2cm]
  {\normalsize Durée : 30 minutes \quad -- \quad Calculatrice autorisée} \\[0.1cm]
  {\normalsize Le barème est donné à titre indicatif.}
\end{center}

\vspace{0.2cm}

\noindent\textbf{Nom :} \dotfill \quad \textbf{Prénom :} \dotfill \quad \textbf{Classe :} \dotfill

\vspace{0.3cm}

\noindent\fbox{\parbox{\dimexpr\textwidth-2\fboxsep-2\fboxrule}{%
\textbf{Rappel :} Soit $f$ une fonction périodique de période $T$. La valeur moyenne de $f$ est :
\[
  a_0 = \frac{1}{T} \int_0^{T} f(t)\, dt
\]
}}

\vspace{0.3cm}

%% ================================================================
%% EXERCICE 1 — 6 points
%% ================================================================
\begin{tcolorbox}[exobox={Exercice 1 \hfill \normalfont\textit{(6 points)}}]
  On considère le signal périodique $f$ de période $T = 4$ défini sur $[0\,;\,4[$ par :
  \[
    f(t) = \begin{cases} 5 & \text{si } 0 \leq t < 2 \\ -3 & \text{si } 2 \leq t < 4 \end{cases}
  \]
\end{tcolorbox}

\begin{center}
  \begin{tikzpicture}
    \begin{axis}[
      axis lines=middle,
      xlabel={$t$},
      ylabel={$f(t)$},
      xmin=-4.5, xmax=8.5,
      ymin=-4.5, ymax=6.5,
      xtick={-4,-2,0,2,4,6,8},
      ytick={-3,-2,-1,0,1,2,3,4,5},
      width=13cm, height=4.8cm,
      grid=major,
      grid style={dashed, gray!30},
    ]
    % Période [-4, 0[
    \addplot[blue, ultra thick, domain=-4:-2] {5};
    \addplot[blue, ultra thick, domain=-2:0] {-3};
    % Période [0, 4[
    \addplot[blue, ultra thick, domain=0:2] {5};
    \addplot[blue, ultra thick, domain=2:4] {-3};
    % Période [4, 8[
    \addplot[blue, ultra thick, domain=4:6] {5};
    \addplot[blue, ultra thick, domain=6:8] {-3};
    % Discontinuités
    \addplot[blue, dashed, thin] coordinates {(-2,5)(-2,-3)};
    \addplot[blue, dashed, thin] coordinates {(2,5)(2,-3)};
    \addplot[blue, dashed, thin] coordinates {(6,5)(6,-3)};
    \end{axis}
  \end{tikzpicture}
\end{center}

\begin{enumerate}[label=\textbf{\arabic*.}]
  \item Donner la période $T$ et les bornes d'intégration. \hfill \textit{(1 pt)}
  \begin{tcolorbox}[reponsebox, height=1.2cm]\end{tcolorbox}

  \item Écrire l'intégrale de $a_0$ en la décomposant sur les deux intervalles. \hfill \textit{(2 pts)}
  \begin{tcolorbox}[reponsebox, height=2cm]\end{tcolorbox}

  \item Calculer $a_0$. \hfill \textit{(2 pts)}
  \begin{tcolorbox}[reponsebox, height=2.5cm]\end{tcolorbox}

  \item Tracer la droite $y = a_0$ sur le graphique. Ce résultat est-il cohérent ? \hfill \textit{(1 pt)}
  \begin{tcolorbox}[reponsebox, height=1.2cm]\end{tcolorbox}
\end{enumerate}

\vspace{0.3cm}

%% ================================================================
%% EXERCICE 2 — 7 points
%% ================================================================
\begin{tcolorbox}[exobox={Exercice 2 \hfill \normalfont\textit{(7 points)}}]
  On considère le signal périodique $g$ de période $T = \pi$ défini sur $[0\,;\,\pi[$ par :
  \[
    g(t) = \begin{cases} \sin(t) & \text{si } 0 \leq t < \dfrac{\pi}{2} \\[6pt] 1 & \text{si } \dfrac{\pi}{2} \leq t < \pi \end{cases}
  \]
\end{tcolorbox}

\begin{center}
  \begin{tikzpicture}
    \begin{axis}[
      axis lines=middle,
      xlabel={$t$},
      ylabel={$g(t)$},
      xmin=-3.5, xmax=7,
      ymin=-0.3, ymax=1.5,
      xtick={-3.14159,-1.5708,0,1.5708,3.14159,4.7124,6.28318},
      xticklabels={$-\pi$,$-\frac{\pi}{2}$,$0$,$\frac{\pi}{2}$,$\pi$,$\frac{3\pi}{2}$,$2\pi$},
      ytick={0,0.5,1},
      width=13cm, height=5cm,
      grid=major,
      grid style={dashed, gray!30},
      samples=100,
    ]
    % Période [-pi, 0[
    \addplot[red, ultra thick, domain=-3.14159:-1.5708] {sin(deg(x+3.14159))};
    \addplot[red, ultra thick, domain=-1.5708:0] {1};
    % Période [0, pi[
    \addplot[red, ultra thick, domain=0:1.5708] {sin(deg(x))};
    \addplot[red, ultra thick, domain=1.5708:3.14159] {1};
    % Période [pi, 2pi[
    \addplot[red, ultra thick, domain=3.14159:4.7124] {sin(deg(x-3.14159))};
    \addplot[red, ultra thick, domain=4.7124:6.28318] {1};
    % Discontinuités
    \addplot[red, dashed, thin] coordinates {(0,1)(0,0)};
    \addplot[red, dashed, thin] coordinates {(3.14159,1)(3.14159,0)};
    \addplot[red, dashed, thin] coordinates {(6.28318,1)(6.28318,0)};
    \end{axis}
  \end{tikzpicture}
\end{center}

\begin{enumerate}[label=\textbf{\arabic*.}]
  \item Écrire l'intégrale de $a_0$ en la décomposant sur les deux intervalles. \hfill \textit{(2 pts)}
  \begin{tcolorbox}[reponsebox, height=2cm]\end{tcolorbox}

  \item Calculer $\displaystyle\int_0^{\pi/2} \sin(t)\,dt$. \hfill \textit{(2 pts)}

  \textit{Rappel :} une primitive de $\sin(t)$ est $-\cos(t)$.
  \begin{tcolorbox}[reponsebox, height=2.5cm]\end{tcolorbox}

  \item En déduire la valeur exacte de $a_0$. \hfill \textit{(2 pts)}
  \begin{tcolorbox}[reponsebox, height=2.5cm]\end{tcolorbox}

  \item Donner une valeur approchée de $a_0$ à $10^{-2}$ près. \hfill \textit{(1 pt)}
  \begin{tcolorbox}[reponsebox, height=1.2cm]\end{tcolorbox}
\end{enumerate}

\vspace{0.3cm}

%% ================================================================
%% EXERCICE 3 — 7 points
%% ================================================================
\begin{tcolorbox}[exobox={Exercice 3 \hfill \normalfont\textit{(7 points)}}]
  Un capteur relève une tension périodique $u$ de période $T = 3$\,ms définie sur $[0\,;\,3[$ (en ms) par :
  \[
    u(t) = \begin{cases} 4t & \text{si } 0 \leq t < 1 \\ 4 & \text{si } 1 \leq t < 2 \\ 0 & \text{si } 2 \leq t < 3 \end{cases}
  \]
  La tension $u(t)$ est exprimée en volts.
\end{tcolorbox}

\begin{enumerate}[label=\textbf{\arabic*.}]
  \item Tracer le signal $u$ sur deux périodes dans le repère ci-dessous. \hfill \textit{(2 pts)}

  \begin{center}
    \begin{tikzpicture}
      \begin{axis}[
        axis lines=middle,
        xlabel={$t$ (ms)},
        ylabel={$u(t)$ (V)},
        xmin=-0.5, xmax=6.5,
        ymin=-0.5, ymax=5.5,
        xtick={0,1,2,3,4,5,6},
        ytick={0,1,2,3,4,5},
        width=13cm, height=5.5cm,
        grid=major,
        grid style={dashed, gray!30},
      ]
      \end{axis}
    \end{tikzpicture}
  \end{center}

  \item Décomposer l'intégrale de $a_0$ sur les trois intervalles. \hfill \textit{(1,5 pts)}
  \begin{tcolorbox}[reponsebox, height=2cm]\end{tcolorbox}

  \item Calculer $a_0$. \hfill \textit{(2,5 pts)}
  \begin{tcolorbox}[reponsebox, height=3.5cm]\end{tcolorbox}

  \item Quelle grandeur physique représente $a_0$ ? Donner sa valeur avec son unité. \hfill \textit{(1 pt)}
  \begin{tcolorbox}[reponsebox, height=1.5cm]\end{tcolorbox}
\end{enumerate}

\end{document}
